\documentclass[12pt, a4paper]{article}

% ── Packages ──────────────────────────────────────────────────────────────────
\usepackage[T1]{fontenc}
\usepackage[utf8]{inputenc}
\usepackage[english]{babel}
\usepackage{lmodern}
\usepackage{microtype}
\usepackage{amsmath, amssymb, amsthm}
\usepackage{bm}
\usepackage{booktabs}
\usepackage{graphicx}
\usepackage[dvipsnames,svgnames,x11names]{xcolor}
\usepackage{hyperref}
\usepackage{geometry}
\usepackage{setspace}
\usepackage{natbib}
\usepackage{caption}
\usepackage{subcaption}

\geometry{margin=2.5cm}
\onehalfspacing

\hypersetup{
    colorlinks=true,
    linkcolor=blue!60!black,
    citecolor=blue!60!black,
    urlcolor=blue!60!black
}

% ── Theorem environments ───────────────────────────────────────────────────────
\newtheorem{remark}{Remark}
\newtheorem{proposition}{Proposition}

% ── Notation macros ───────────────────────────────────────────────────────────
\newcommand{\E}{\mathbb{E}}
\newcommand{\Var}{\mathrm{Var}}
\newcommand{\tr}{\mathrm{tr}}
\newcommand{\diag}{\mathrm{diag}}
\newcommand{\bs}[1]{\boldsymbol{#1}}
\newcommand{\eye}{\mathbf{I}}

% ══════════════════════════════════════════════════════════════════════════════
\begin{document}

\begin{titlepage}
    \centering
    \vspace*{2cm}
    {\LARGE \textbf{VAR with Shifting Endpoint:} \\[0.4em]
    \textbf{Long-Run Anchoring via Focus Survey Expectations}} \\[1.5cm]
    {\large A Replication of BCB Estudo Especial 19/2018} \\[2cm]
    {\large \today}
    \vfill
    {\small Based on Kozicki \& Tinsley (2012), Clark \& West (2007), and BCB (2018)}
\end{titlepage}

\tableofcontents
\newpage

% ══════════════════════════════════════════════════════════════════════════════
\section{Introduction}

Standard Vector Autoregressive (VAR) models assume that the unconditional mean
of the system is fixed. For the Brazilian inflation process, this assumption is
problematic over a sample spanning the mid-2000s to the present: the implicit
inflation target has declined from around 8\% to 3.5\% per annum, and long-run
private-sector expectations---as measured by the Focus survey---have shifted
substantially over time.

The \textbf{VAR with Shifting Endpoint} (VAR-SE), introduced by
\citet{kozicki2012} and applied to Brazil by the Banco Central do Brasil (BCB)
in Estudo Especial 19/2018, addresses this by allowing the long-run mean of
inflation---the \emph{endpoint}---to be a latent, time-varying random walk
$\mu_t$. The endpoint is anchored to Focus survey expectations via a
state-space framework and a Kalman filter. At long horizons the conditional
forecast converges to $\mu_t$ rather than the sample mean, producing
forward-looking paths that respect current market expectations.

The model estimated here uses three endogenous variables:
\begin{itemize}
    \item $\pi_t$: monthly IPCA inflation (\%/month);
    \item $\Delta e_t$: monthly BRL/USD exchange-rate change (\%/month);
    \item $r_t^{\mathrm{real}}$: ex-post real interest rate (\%/month).
\end{itemize}
The shifting endpoint $\mu_t$ applies to the inflation equation only.

% ══════════════════════════════════════════════════════════════════════════════
\section{The Unrestricted VAR}

Let $\bs{y}_t = (\pi_t,\, \Delta e_t,\, r_t^{\mathrm{real}})'$ be the
$k\times 1$ vector of endogenous variables ($k=3$). The benchmark
unrestricted VAR($p$) is
\begin{equation}
    \bs{y}_t \;=\; \bs{c} + \sum_{j=1}^{p} \bs{\Phi}_j \bs{y}_{t-j} + \bs{\varepsilon}_t,
    \qquad \bs{\varepsilon}_t \sim \mathcal{N}(\bs{0},\,\bs{\Sigma}),
    \label{eq:var}
\end{equation}
where $\bs{c}\in\mathbb{R}^k$ is the intercept vector and
$\bs{\Phi}_j\in\mathbb{R}^{k\times k}$ are coefficient matrices. Equation
\eqref{eq:var} is estimated equation-by-equation by OLS.

\subsection{Long-Run Mean}

The unconditional mean implied by the VAR is
\begin{equation}
    \bar{\bs{y}} \;=\; \left(\eye_k - \sum_{j=1}^{p}\bs{\Phi}_j\right)^{-1}\bs{c}.
    \label{eq:lr_mean}
\end{equation}
After stability enforcement (see Section~\ref{sec:stability}), the intercept
is corrected to preserve $\bar{\bs{y}}$ at the sample mean $\widehat{\bs{y}}$:
\begin{equation}
    \bs{c} \;\leftarrow\; \left(\eye_k - \sum_{j=1}^{p}\hat{\bs{\Phi}}_j\right)\widehat{\bs{y}}.
    \label{eq:intercept_correction}
\end{equation}

\subsection{Stability Enforcement}
\label{sec:stability}

Let $\rho$ denote the spectral radius of the VAR companion matrix (defined
below). If $\rho \geq \rho_{\max}$ (we use $\rho_{\max}=0.98$), all coefficient
matrices are scaled uniformly:
\begin{equation}
    \hat{\bs{\Phi}}_j \;\leftarrow\; \frac{\rho_{\max}}{\rho}\,\hat{\bs{\Phi}}_j,
    \qquad j=1,\ldots,p.
\end{equation}
The intercept is then recomputed via \eqref{eq:intercept_correction} to
maintain the unconditional mean.

% ══════════════════════════════════════════════════════════════════════════════
\section{VAR with Shifting Endpoint}
\label{sec:varse}

\subsection{The Endpoint Process}

The VAR-SE replaces the fixed intercept $\bs{c}$ in the inflation equation
with a time-varying endpoint $\mu_t$. The endpoint follows a random walk:
\begin{equation}
    \mu_t \;=\; \mu_{t-1} + \eta_t, \qquad \eta_t \sim \mathcal{N}(0,\,\sigma^2_\mu).
    \label{eq:mu_rw}
\end{equation}
The VAR-SE can be written as
\begin{equation}
    \pi_t \;=\; (1-\phi_{11,1}-\cdots-\phi_{11,p})\,\mu_t
                + \sum_{j=1}^{p}\sum_{l=1}^{k}\phi_{1l,j}\,y_{l,t-j} + \varepsilon_{1,t},
    \label{eq:varse_pi}
\end{equation}
where $\phi_{1l,j}$ is element $(1,l)$ of $\bs{\Phi}_j$. The remaining
equations ($\Delta e_t$, $r_t^{\mathrm{real}}$) retain fixed intercepts.

The key structural property of \eqref{eq:varse_pi} is that as
$h\to\infty$ the $h$-step conditional forecast satisfies
$\E_t[\pi_{t+h}] \to \mu_t$. The endpoint is therefore the long-run
inflation forecast implied by the model.

\subsection{Deviation-from-Endpoint Form}

Define $\tilde{\bs{y}}_t = \bs{y}_t - \bs{e}_1\mu_t$ where
$\bs{e}_1 = (1,0,0)'$. In deviations the inflation equation becomes
\begin{equation}
    \tilde{\pi}_t \;=\; \sum_{j=1}^{p}\bs{\Phi}_{1\cdot,j}\tilde{\bs{y}}_{t-j}
                        + \varepsilon_{1,t},
\end{equation}
with no constant. This is the \emph{deviation-from-endpoint} representation
used to construct the companion matrix (Section~\ref{sec:companion}).

% ══════════════════════════════════════════════════════════════════════════════
\section{State-Space Representation}
\label{sec:statespace}

\subsection{Companion Matrix}
\label{sec:companion}

Stack the endogenous lags into the $(kp+2)\times 1$ state vector
\begin{equation}
    \bs{X}_t \;=\; \left(\bs{y}_t',\;\bs{y}_{t-1}',\;\ldots,\;\bs{y}_{t-p+1}',\;1,\;\mu_t\right)',
    \label{eq:state}
\end{equation}
where the constant 1 occupies position $kp+1$ (0-indexed: $kp$) and $\mu_t$
occupies position $kp+2$ (0-indexed: $kp+1$). The transition equation is
\begin{equation}
    \bs{X}_t \;=\; \bs{A}\,\bs{X}_{t-1} + \bs{w}_t,
    \label{eq:transition}
\end{equation}
where $\bs{w}_t$ has covariance $\bs{Q}$ with $Q_{\mu\mu}=\sigma^2_\mu$ and
zeros elsewhere.

The companion matrix $\bs{A}$ has the standard VAR block structure, augmented
by the endpoint column. For the VAR-SE, the first row satisfies
\begin{equation}
    A_{1,\mu} \;=\; 1 - \sum_{j=1}^{p}\Phi_{11,j},
    \label{eq:varse_col}
\end{equation}
which encodes the deviation-from-endpoint property. The unrestricted VAR uses
$A_{1,\mu}=1$, recovering the OLS unconditional mean.

Explicitly, $\bs{A}$ is:
\begin{equation}
    \bs{A} \;=\;
    \begin{pmatrix}
        \bs{\Phi}_1 & \bs{\Phi}_2 & \cdots & \bs{\Phi}_p & \bs{c}^* & \bs{a}_\mu \\
        \eye_k      & \bs{0}      & \cdots & \bs{0}      & \bs{0}   & \bs{0}     \\
        \bs{0}      & \eye_k      & \cdots & \bs{0}      & \bs{0}   & \bs{0}     \\
        \vdots      &             & \ddots & \vdots      & \vdots   & \vdots     \\
        \bs{0}      & \bs{0}      & \cdots & \eye_k      & \bs{0}   & \bs{0}     \\
        \bs{0}      & \bs{0}      & \cdots & \bs{0}      & 1        & 0          \\
        \bs{0}      & \bs{0}      & \cdots & \bs{0}      & 0        & 1
    \end{pmatrix},
    \label{eq:companion}
\end{equation}
where $\bs{a}_\mu$ is a $k\times 1$ column with $A_{1,\mu}$ in the first
position and zeros elsewhere, and $\bs{c}^*$ contains the fixed-intercept
contributions from $\bs{c}$ for equations 2 through $k$ (zero for the
inflation equation, which uses $\mu_t$).

\subsection{The Survey Constraint}

The Focus survey provides the median 4-year-ahead annual IPCA expectation,
$f_t$ (\%/a.a.). This corresponds to the 12-month accumulated inflation from
month $t+37$ to $t+48$. The linear constraint linking $f_t$ to the state
is
\begin{equation}
    f_t \;=\; \bs{c}_{\mathrm{sv}}'\bs{X}_t + u_t,
    \qquad u_t \sim \mathcal{N}(0,\,\sigma^2_{\mathrm{sv}}),
    \label{eq:survey}
\end{equation}
where $\bs{c}_{\mathrm{sv}}' = \bs{e}_1'\sum_{h=37}^{48}\bs{A}^h$ and
$\sigma^2_{\mathrm{sv}}$ is the survey measurement-noise variance.

The survey row $\bs{c}_{\mathrm{sv}}'$ is constructed from the VAR-SE
companion matrix $\bs{A}$ by summing the propagation matrices over the
relevant horizon window. Its sensitivity to $\mu_t$ (the element
$c_{\mathrm{sv},\mu}$) is approximately 12, reflecting the accumulation
of 12 monthly rates each of which responds one-for-one to $\mu_t$ at
long horizons.

\subsection{How the Survey Identifies \texorpdfstring{$\mu_t$}{mu\_t}}
\label{sec:survey_id}

Expanding the dot product $\bs{c}_{\mathrm{sv}}'\bs{X}_t$ using the state
vector definition \eqref{eq:state}:
\begin{equation}
    \bs{c}_{\mathrm{sv}}'\bs{X}_t
    \;=\;
    \underbrace{\bs{c}_{\mathrm{sv}}^{[y]}\cdot \bs{y}_t^{(\mathrm{stack})}}_{\text{current and lagged }y_t}
    \;+\;
    \underbrace{c_{\mathrm{sv},\bar{x}}}_{\approx\,0}
    \;+\;
    \underbrace{c_{\mathrm{sv},\mu}}_{\approx\,12}
    \cdot\mu_t,
    \label{eq:survey_expand}
\end{equation}
where $c_{\mathrm{sv},\mu} = \bs{e}_1'\sum_{h=37}^{48}\bs{A}^h\,\bs{e}_\mu$
is the sensitivity of the 37--48 month cumulative forecast to $\mu_t$.

\paragraph{Why $c_{\mathrm{sv},\mu}\approx 12$.}
Under stationarity, $\bs{A}^h \to \bs{e}_\mu\bs{e}_\mu'$ (projection onto
the unit-eigenvalue direction) as $h\to\infty$. At $h \gtrsim 37$ months
this limit is already a good approximation, so
\begin{equation}
    \bs{e}_1'\bs{A}^h\bs{e}_\mu \;\approx\; 1,
    \qquad h \geq 37,
    \label{eq:attenuation}
\end{equation}
meaning each monthly forecast from $h=37$ to $h=48$ responds one-for-one to
$\mu_t$. Summing over 12 months gives $c_{\mathrm{sv},\mu}\approx 12$.

\paragraph{Why $\bs{c}_{\mathrm{sv}}^{[y]}\approx \bs{0}$.}
The contribution of the current VAR state $\bs{y}_t$ to the forecast at
horizon $h$ decays as $\rho^h$, where $\rho < 1$ is the spectral radius of
the stationary block of $\bs{A}$. With $\rho \leq 0.96$ and $h \geq 37$,
we have $\rho^{37} \leq 0.22$; the 12-period sum
$\sum_{h=37}^{48}\rho^h \leq 2.4$ is small relative to $c_{\mathrm{sv},\mu}=12$.
Concretely, the survey row is almost entirely loaded on $\mu_t$, not on
$\bs{y}_t$.

\paragraph{Consequence.}
Rearranging \eqref{eq:survey} and \eqref{eq:survey_expand}:
\begin{equation}
    \mu_t \;\approx\;
    \frac{f_t \;-\; \bs{c}_{\mathrm{sv}}^{[y]}\cdot\bs{y}_t^{(\mathrm{stack})}}{c_{\mathrm{sv},\mu}}
    \;-\; \frac{u_t}{c_{\mathrm{sv},\mu}}.
    \label{eq:mu_from_focus}
\end{equation}
The endpoint $\mu_t$ is therefore approximately equal to the Focus
48m expectation divided by 12 (to convert annual to monthly), corrected
for the small residual contribution of the current VAR state at long
horizons. The Kalman filter implements this identification softly,
weighting the survey signal against the prior variance
$P_{t|t-1,\mu\mu}$ via the gain
$K_\mu = P_{t|t-1,\mu\mu}\,c_{\mathrm{sv},\mu} / s^2_t$.

% ══════════════════════════════════════════════════════════════════════════════
\section{Kalman Filter}
\label{sec:kalman}

\subsection{Key Observation}

The endogenous variables $\bs{y}_t$ are measured without error: they are
observed directly. Only the latent endpoint $\mu_t$ is uncertain.
Accordingly, the filter treats $\bs{y}_t$ as an exact observation and
performs a \emph{scalar} Kalman update using only the Focus survey
innovation. This prevents the numerical degeneracy that arises when $k$
near-zero-noise observation equations are combined with a rank-1 covariance
matrix for $\bs{X}_t$.

\subsection{Filter Recursion}

Given $\bs{X}_{t-1|t-1}$ and $\bs{P}_{t-1|t-1} = \Var(\bs{X}_{t-1}|\mathcal{F}_{t-1})$,
the filter iterates:

\medskip
\noindent\textbf{Step 1 — Prediction:}
\begin{align}
    \bs{X}_{t|t-1} &\;=\; \bs{A}\,\bs{X}_{t-1|t-1}, \label{eq:predict_x} \\
    \bs{P}_{t|t-1} &\;=\; \bs{A}\,\bs{P}_{t-1|t-1}\,\bs{A}' + \bs{Q}. \label{eq:predict_p}
\end{align}

\noindent\textbf{Step 2 — Scalar survey update:}
\begin{align}
    v_t   &\;=\; f_t - \bs{c}_{\mathrm{sv}}'\bs{X}_{t|t-1}, \label{eq:innov} \\
    s^2_t &\;=\; \bs{c}_{\mathrm{sv}}'\bs{P}_{t|t-1}\bs{c}_{\mathrm{sv}} + \sigma^2_{\mathrm{sv}}, \label{eq:s2} \\
    \bs{K}_t &\;=\; \frac{\bs{P}_{t|t-1}\,\bs{c}_{\mathrm{sv}}}{s^2_t}, \label{eq:gain} \\
    \bs{X}_{t|t} &\;=\; \bs{X}_{t|t-1} + \bs{K}_t\,v_t, \label{eq:update_x} \\
    \bs{P}_{t|t} &\;=\; \bs{P}_{t|t-1} - \bs{K}_t\bs{K}_t'\,s^2_t. \label{eq:update_p}
\end{align}

\noindent\textbf{Step 3 — Exact observation override:}

Since $\bs{y}_t$ is observed without error, the VAR state components are
overridden from data:
\begin{equation}
    X_{t|t,l} \;=\; y_{t-\lfloor l/k \rfloor, (l \bmod k)+1},
    \quad l = 0,\ldots,kp-1,
    \qquad X_{t|t,kp} = 1.
    \label{eq:override}
\end{equation}
Correspondingly, all rows and columns of $\bs{P}_{t|t}$ associated with
directly-observed components are set to zero, leaving $P_{t|t,\mu\mu}$ as
the sole uncertain element.

\medskip
\noindent\textbf{Log-likelihood contribution:}
\begin{equation}
    \ell_t \;=\; -\tfrac{1}{2}\left(\log s^2_t + \frac{v_t^2}{s^2_t}\right).
    \label{eq:ll}
\end{equation}
The total log-likelihood is $\mathcal{L} = \sum_{t=p}^{T-1}\ell_t$.

\subsection{Initialization}

The state is initialized at $t=p$ (the first fully-observed period of the
VAR sample):
\begin{equation}
    \bs{X}_{p|p-1} \;=\;
    \left(\bs{y}_p',\;\bs{y}_{p-1}',\;\ldots,\;\bs{y}_1',\;1,\;\hat{c}_1\right)',
\end{equation}
where $\hat{c}_1$ is the OLS intercept for the inflation equation (a
diffuse, data-consistent starting value for $\mu_0$). The initial covariance
is $\bs{P}_{p|p-1} = \diag(0,\ldots,0,0,P_0)$ with $P_0=10$ (diffuse prior
on $\mu_0$) and zeros for all directly-observed components.

% ══════════════════════════════════════════════════════════════════════════════
\section{Estimation}
\label{sec:estimation}

\subsection{OLS Step}

The coefficient matrices $\{\bs{\Phi}_j\}$ and $\bs{c}$ are estimated by
OLS on the unrestricted VAR \eqref{eq:var}. The OLS estimates are then
used unchanged in the VAR-SE companion matrix; only $\mu_t$ is re-estimated
via the Kalman filter.

This two-step approach follows BCB (2018): the VAR dynamics are identified
from the full sample, and the endpoint captures the slowly-drifting
long-run mean not adequately described by a fixed intercept.

\subsection{Calibration of \texorpdfstring{$\sigma^2_{\mathrm{sv}}$}{sigma2\_survey}}

The measurement-noise variance $\sigma^2_{\mathrm{sv}}$ governs how
tightly $\mu_t$ is anchored to Focus. It cannot be identified by MLE
in the scalar-Kalman setup: the log-likelihood is monotonically increasing
as $\sigma^2_{\mathrm{sv}}\to 0$ (perfect anchoring), rendering the
problem degenerate.

Following the Kozicki--Tinsley convention, $\sigma^2_{\mathrm{sv}}$ is
calibrated externally. We set
$\sigma^2_{\mathrm{sv}} = 0.25\;(\%\text{ a.a.})^2$, corresponding to a
standard deviation of 0.5 pp a.a., consistent with the typical
cross-sectional dispersion of Focus respondents at the 4-year horizon.

\begin{remark}
    The calibration implies that a 95\% credible interval for the survey
    observation is $\pm 1$ pp around the median forecast. Tighter
    calibration ($\sigma^2_{\mathrm{sv}}\downarrow$) produces $\mu_t$
    paths that more closely mirror Focus; looser calibration
    ($\sigma^2_{\mathrm{sv}}\uparrow$) produces smoother, more inertial
    paths.
\end{remark}

\subsection{MLE for \texorpdfstring{$\sigma^2_\mu$}{sigma2\_mu}}

Conditional on $\sigma^2_{\mathrm{sv}}$, the drift variance $\sigma^2_\mu$
is estimated by maximising the Kalman log-likelihood:
\begin{equation}
    \hat{\sigma}^2_\mu \;=\; \arg\max_{\sigma^2_\mu \in [10^{-8},\,10^{-2}]}
    \mathcal{L}\!\left(\sigma^2_\mu\,;\,\sigma^2_{\mathrm{sv}}\right).
    \label{eq:mle}
\end{equation}
The one-dimensional optimisation uses Brent's bounded method on the
log-scale $\log\sigma^2_\mu$ for numerical stability.

% ══════════════════════════════════════════════════════════════════════════════
\section{Forecasting}

\subsection{Initial Conditions}

At the forecast origin $T$, the initial state for the VAR-SE is
\begin{equation}
    \bs{X}_T \;=\; \left(\bs{y}_T',\;\bs{y}_{T-1}',\;\ldots,\;\bs{y}_{T-p+1}',\;1,\;\hat\mu_T\right)',
\end{equation}
where $\hat\mu_T = X_{T|T,\mu}$ is the filtered endpoint. For the
unrestricted VAR, $\mu_T$ is replaced by $\hat{c}_1$ (the OLS intercept for
the inflation equation).

\subsection{$h$-Step Forecasts}

The $h$-step ahead point forecast is
\begin{equation}
    \widehat{\bs{y}}_{T+h} \;=\; \bs{J}\,\bs{A}^h\,\bs{X}_T,
\end{equation}
where $\bs{J} = [\eye_k\;\bs{0}_{k\times 2}]$ selects the first $k$
components of the state. The 12-month accumulated inflation forecast from
horizon $h_1$ to $h_2$ is
\begin{equation}
    \widehat{\Pi}_{T}^{h_1:h_2} \;=\; \sum_{h=h_1}^{h_2}\widehat{\pi}_{T+h}.
    \label{eq:accum}
\end{equation}

\subsection{Long-Run Convergence}

Under the VAR-SE companion matrix, $\bs{A}^h \to \bs{e}_\mu\bs{e}_\mu'$ as
$h\to\infty$ (where $\bs{e}_\mu$ is the unit vector selecting $\mu$), so
\begin{equation}
    \widehat{\pi}_{T+h} \;\to\; \hat\mu_T \quad \text{as } h\to\infty.
\end{equation}
Consequently the 48-month accumulated forecast satisfies
$\widehat{\Pi}_T^{1:48}/12 \approx \hat\mu_T$, anchored to the Focus
survey median.

For the unrestricted VAR, $\widehat{\pi}_{T+h}\to\bar\pi$ (the sample mean
of inflation, $\approx 5.9\%$ a.a. over the 2001--2026 sample), which
departs significantly from the current Focus expectation of 3.5\% a.a.

% ══════════════════════════════════════════════════════════════════════════════
\section{Out-of-Sample Evaluation}
\label{sec:oos}

\subsection{Pseudo-Real-Time Exercise}

Forecasting accuracy is assessed by an expanding-window pseudo-real-time
exercise starting from 2013:06. At each period $t$, the VAR is re-estimated
on $\{\bs{y}_1,\ldots,\bs{y}_t\}$, the Kalman filter is run with fixed
$(\hat\sigma^2_\mu,\hat\sigma^2_{\mathrm{sv}})$ from the full-sample
estimation, and $h$-step-ahead forecasts are generated for
$h\in\{1,3,6,12,24,36\}$ months.

\subsection{Loss Function}

The benchmark loss is Mean Squared Error (MSE):
\begin{equation}
    \mathrm{MSE}^{(h)} \;=\; \frac{1}{N_h}\sum_{t}\left(
        \widehat{\Pi}_{t}^{1:h} - \Pi_{t}^{1:h}\right)^2,
\end{equation}
where $\Pi_t^{1:h}$ is the realised 12-month accumulated inflation computed
from subsequent observations and $N_h$ is the number of forecast origins with
available realisations.

\subsection{Statistical Tests}

Two tests are used:

\medskip
\noindent\textbf{Diebold--Mariano (DM)} \citep{dm1995}: tests the null of
equal unconditional expected loss between the VAR-SE and a benchmark:
\begin{equation}
    H_0:\; \E\!\left[d_{t+h}\right] = 0,
    \qquad d_{t+h} = e^2_{\text{bench},t+h} - e^2_{\text{VARSE},t+h}.
\end{equation}
The DM statistic uses HAC-robust variance estimation (Newey--West with
$h-1$ lags). A positive test statistic (small p-value) means VAR-SE
has lower expected loss.

\medskip
\noindent\textbf{Clark--West (CW)} \citep{clarkewest2007}: corrects for
the upward bias in the VAR-SE MSE that arises mechanically when the null
model is nested. The adjusted loss differential is
\begin{equation}
    \tilde{d}_{t+h} \;=\; d_{t+h} + \left(\widehat{\Pi}^{\text{bench}}_{t+h} - \widehat{\Pi}^{\text{VARSE}}_{t+h}\right)^2,
\end{equation}
and the CW statistic is the t-ratio of $\bar{\tilde{d}}$ under the null
$H_0: \E[\tilde{d}_{t+h}]\leq 0$. This is the recommended test when the
models are nested (the unrestricted VAR is nested in the VAR-SE when
$\sigma^2_\mu\to 0$).

% ══════════════════════════════════════════════════════════════════════════════
\section{Empirical Results}
\label{sec:results}

\subsection{Full-Sample Estimation}

Table~\ref{tab:params} reports the estimated parameters for the VAR(2)-SE
using the sample 2001:11--2026:02 ($T=292$ observations, $k=3$, $p=2$).

\begin{table}[h!]
\centering
\caption{Full-sample parameter estimates, VAR(2)-SE}
\label{tab:params}
\begin{tabular}{lcc}
\toprule
Parameter & Symbol & Estimate \\
\midrule
Spectral radius (VAR companion) & $\rho$ & 0.9598 \\
Endpoint drift variance         & $\hat\sigma^2_\mu$ & $3.0\times10^{-5}$ \\
Survey noise variance (calibrated) & $\sigma^2_{\mathrm{sv}}$ & 0.25 \\
Kalman log-likelihood           & $\mathcal{L}$ & 144.51 \\
\midrule
OLS mean: $\pi$     & $\bar\pi$      & 0.492\%/m (5.90\% a.a.) \\
OLS mean: $\Delta e$& $\bar{\Delta e}$ & 0.340\%/m (4.07\% a.a.) \\
OLS mean: $r^r$     & $\bar r^r$     & 0.441\%/m (5.29\% a.a.) \\
Filtered endpoint (last)& $\hat\mu_T$ & 0.273\%/m (3.27\% a.a.) \\
Focus 48m (last)    & $f_T$          & 3.50\% a.a. \\
\bottomrule
\end{tabular}
\end{table}

The filtered endpoint $\hat\mu_T = 0.273$\%/m lies between the OLS
sample mean (0.492\%/m) and Focus/12 (0.292\%/m), reflecting the
Kalman filter's balancing of prior inertia against the survey signal.
The small $\hat\sigma^2_\mu = 3\times10^{-5}$ implies the endpoint
moves slowly (annualised standard deviation
$\approx\sqrt{12\times3\times10^{-5}}\approx 0.02$\%/m $\approx 0.24$
pp a.a.), consistent with a gradual de-anchoring/re-anchoring process.

\subsection{Long-Run Forecasts}

Table~\ref{tab:lr} compares the 48-month accumulated inflation forecasts
across models. The unrestricted VAR and ARMA converge to the sample mean
($\approx 5.9$\% a.a.), well above the Focus survey (3.50\% a.a.).
The VAR-SE converges to $\hat\mu_T$ (3.55\% a.a.), within 5 basis points
of Focus.

\begin{table}[h!]
\centering
\caption{48-month accumulated inflation forecast (annualised, \% a.a.)}
\label{tab:lr}
\begin{tabular}{lc}
\toprule
Model & Forecast \\
\midrule
ARMA(4,3)   & 5.90 \\
VAR(2)      & 5.93 \\
VAR(2)-SE   & 3.55 \\
Focus (observed) & 3.50 \\
\bottomrule
\end{tabular}
\end{table}

\subsection{Out-of-Sample MSE}

Table~\ref{tab:oos} reports MSE ($\times10^4$) and p-values for the
Diebold--Mariano and Clark--West tests. Entries for DM and CW are
p-values of the null that VAR-SE is no better than the benchmark;
asterisks indicate 10\% (*), 5\% (**), and 1\% (***) significance.

\begin{table}[h!]
\centering
\caption{Out-of-sample evaluation, 2013:06 onward}
\label{tab:oos}
\begin{tabular}{rcccccccc}
\toprule
$h$ & MSE$_{\text{ARMA}}$ & MSE$_{\text{VAR}}$ & MSE$_{\text{VARSE}}$
    & $N$ & $p_{\text{DM,A}}$ & $p_{\text{DM,V}}$
    & $p_{\text{CW,A}}$ & $p_{\text{CW,V}}$ \\
\midrule
 1  & 1107  & 1126  & 1128  & 153 & 0.733 & 0.923 & 0.080*   & 0.261 \\
 3  & 7556  & 7503  & 7521  & 151 & 0.967 & 0.964 & 0.006*** & 0.114 \\
 6  & 24010 & 23950 & 24234 & 148 & 0.957 & 0.900 & 0.010**  & 0.045** \\
12  & 59264 & 64570 & 67614 & 142 & 0.652 & 0.769 & 0.006*** & 0.006*** \\
24  & 186342& 212557& 230574& 130 & 0.527 & 0.702 & 0.000*** & 0.001*** \\
36  & 284017& 323334& 384904& 118 & 0.415 & 0.536 & 0.000*** & 0.000*** \\
\bottomrule
\end{tabular}
\end{table}

Two patterns stand out. First, the raw MSE of VAR-SE is higher than both
benchmarks at all horizons. This reflects the well-known result that
anchoring to a slowly-moving endpoint introduces a downward bias in
near-term forecasts during periods when inflation is above target (as was
the case for much of 2013--2025 in Brazil). Second, after the Clark--West
correction the VAR-SE significantly outperforms the unrestricted VAR at
$h \geq 6$ months. The value of the shifting endpoint is therefore
concentrated at medium and long horizons, consistent with the theoretical
convergence property in Section~6.3.

% ══════════════════════════════════════════════════════════════════════════════
\section{Conclusion}

The VAR with Shifting Endpoint provides a coherent framework for
incorporating long-run survey expectations into a multivariate time-series
model. The Kalman filter jointly updates the endpoint on the basis of the
Focus survey and the VAR dynamics, with the degree of anchoring governed
by the ratio $\sigma^2_\mu / \sigma^2_{\mathrm{sv}}$.

The empirical results for Brazil over 2001--2026 confirm the BCB (2018)
finding: the VAR-SE does not improve short-horizon (1--3 month) point
forecast accuracy, but it meaningfully reduces long-horizon errors and
produces 12-to-48-month paths consistent with market expectations. The
Clark--West test, which accounts for the nesting bias, shows statistically
significant improvements at $h \geq 6$.

For practical use, the model's main output is the filtered endpoint series
$\{\hat\mu_t\}$, which provides a real-time, model-consistent estimate of
where the inflation process is expected to settle in the long run.

% ══════════════════════════════════════════════════════════════════════════════
\bibliographystyle{apalike}
\begin{thebibliography}{9}

\bibitem[BCB(2018)]{bcb2018}
Banco Central do Brasil (2018).
\newblock Modelo de vetor autorregressivo com ancoragem de longo prazo.
\newblock \emph{Estudo Especial} No.\ 19/2018.

\bibitem[Clark and West(2007)]{clarkewest2007}
Clark, T.~E. and West, K.~D. (2007).
\newblock Approximately normal tests for equal predictive accuracy in nested
  models.
\newblock \emph{Journal of Econometrics}, 138(1):291--311.

\bibitem[Diebold and Mariano(1995)]{dm1995}
Diebold, F.~X. and Mariano, R.~S. (1995).
\newblock Comparing predictive accuracy.
\newblock \emph{Journal of Business \& Economic Statistics}, 13(3):253--263.

\bibitem[Kozicki and Tinsley(2012)]{kozicki2012}
Kozicki, S. and Tinsley, P.~A. (2012).
\newblock Moving endpoints and the internal consistency of agents' ex ante
  forecasts.
\newblock \emph{Computational Economics}, 39(1):67--98.

\end{thebibliography}

\end{document}
